%! Systems Involving Quadratics — Unit 3, Lesson 3.6 — Homework (Answer Key)
\documentclass[10pt]{article}
\usepackage{algebra2-article}
\usepackage{algebra2-key}   % pulls in -boxes; do NOT also load -boxes
\newcommand{\TallMath}[1]{$\displaystyle #1\rule[-1.4em]{0pt}{3.2em}$}

\begin{document}

\pageheader{Unit 3, Lesson 3.6}{Homework --- Answer Key}
\namedateperiod

\vspace{0.10in}

\begin{notesbox}{Practice}
Show your work. Every solution is an \textbf{ordered pair} $(x,y)$ --- find both coordinates and check them in
\emph{both} equations.

\begin{enumerate}[leftmargin=1.9em, itemsep=9pt]
  \item \textbf{Is the pair a solution?} Check each against $\ y=x^2+2\ $ and $\ y=2x+1$ (write yes/no).
    \[ (1,3):\ \ans{yes} \qquad\qquad (0,2):\ \ans{no (fails the line: $2(0)+1=1\ne2$)} \]

  \item \textbf{Solve by substitution} (linear--quadratic; real, rational). Write ordered pairs.
    \begin{enumerate}[label=(\alph*), leftmargin=1.9em, itemsep=4pt]
      \item $y=x^2-2x-3,\ \ y=2x-3 \Rightarrow$ \ans{$(0,-3)$ and $(4,5)$}
      \item $y=x^2-4x+3,\ \ y=x-1 \Rightarrow$ \ans{$(1,0)$ and $(4,3)$}
    \end{enumerate}

  \item \textbf{Won't factor} --- use the quadratic formula. Give the $x$-coordinates in simplest radical form.
    \[ y=x^2+1,\ \ y=3x \Rightarrow x=\ans{$\dfrac{3\pm\sqrt{5}}{2}$} \]

  \item \textbf{Quadratic--quadratic.} Set the expressions equal.
    \begin{enumerate}[label=(\alph*), leftmargin=1.9em, itemsep=4pt]
      \item $y=x^2,\ \ y=-x^2+18 \Rightarrow$ \ans{$(-3,9)$ and $(3,9)$}
      \item $y=x^2+2x,\ \ y=x^2-4 \Rightarrow$ \ans{$(-2,0)$ only} \quad\emph{(watch the $x^2$ terms!)}
    \end{enumerate}

  \item \textbf{How many solutions?} Each system collapses to the quadratic shown. Compute $b^2-4ac$ and state how
        many points the graphs share ($0$, $1$, or $2$).
    \[ x^2-4x+4=0:\ \ans{$0\Rightarrow 1$ pt} \qquad x^2-x-6=0:\ \ans{$25\Rightarrow 2$ pts} \qquad x^2+x+3=0:\ \ans{$-11\Rightarrow 0$ pts} \]

  \item \textbf{Find the error.} Priya solves $\ y=x^2,\ \ y=x+6\ $ and correctly gets $x=3$ and $x=-2$. She writes
        ``the solutions are $x=3$ and $x=-2$.'' What did she forget, and what are the correct solutions?
        \ansline{She stopped at the $x$-values --- a solution of a system is an ordered \emph{pair}. Back-substitute into $y=x+6$: the solutions are $(3,9)$ and $(-2,4)$.}
\end{enumerate}
\end{notesbox}

\begin{extensionbox}
\textbf{1. Model (break-even).} A start-up's revenue is $y=-x^2+12x$ and its cost is $y=4x+12$ (hundreds of
dollars, $x$ in hundreds of units). Find both break-even points and the interval where it turns a profit.
\[ x=\ans{$2,\ 6$} \quad\text{points: } \ans{$(2,20),(6,36)$} \quad\text{profit when } \ans{$2<x<6$} \]

\textbf{2. Make it tangent.} For what value of $c$ is the line $y=4x+c$ tangent to $y=x^2$ (exactly one solution)?
\emph{Set equal, then force $b^2-4ac=0$.} \ans{$c=-4$ (then $x^2-4x+4=0$, tangent at $(2,4)$)}
\end{extensionbox}

\begin{spiralbox}
\textbf{Coming next --- Lesson 3.7: Modeling with Quadratics.} You've now solved single quadratics \emph{and}
systems. Next you'll build the quadratic model yourself from a projectile, area, or optimization story --- and read
its \emph{vertex} as the maximum height, maximum area, or best price. The algebra is all from this unit; the new
skill is translating the situation into the model and interpreting the answer.
\end{spiralbox}

\begin{teachernote}
Common slips to look for: \#2 stopping at $x$ without pairing coordinates, or back-substituting into the parabola and
erring; \#3 leaving $\frac{3\pm\sqrt5}{2}$ unsimplified or trying to force factors; \#4(b) is the key
quadratic--quadratic case --- the $x^2$ terms cancel, leaving the \emph{linear} $2x=-4$, so there is exactly
\emph{one} solution (students who expect two will hunt for a phantom second point). \#5 rehearses ``discriminant $=$
intersection count'' (note $D=0$ means \emph{one} point, not zero). \#6 is the lesson's headline misconception ---
credit any answer that names the missing $y$-coordinates and writes $(3,9),(-2,4)$. Extension \#2 reuses the
Lesson 3.5 discriminant as a design dial.
\end{teachernote}

\end{document}
